\section{Method}
\label{sec:method}

\subsection{Problems, instances and bounds}
\label{sec:instances}
\paragraph{Identical parallel machines.} In $P||\Cmax$, $n$ jobs with processing times $p_j$ are assigned to $m$ identical machines to minimise $\Cmax=\max_i\sum_{j\in S_i}p_j$. The classical bound is $\mathrm{LB}=\max(\max_jp_j,\lceil\sum_jp_j/m\rceil)$~\cite{Graham1969}, and $\mathrm{GAP}=100(\Cmax-\mathrm{LB})/\mathrm{LB}$; a tighter, certified bound comes from the linear relaxation of the arc-flow model~\cite{MradSouayah2018} and equals the best makespan found in $75.3\%$ of the synthetic instances (which proves their optimality). The \emph{synthetic} instances have $p_j\sim U[1,100]$ in twelve classes $(n,m)$ from $(20,4)$ to $(300,20)$, with $100$ test instances and $30$ separate tuning instances per class. With integer times one unit of makespan is $0.12$--$0.78\%$ of it (median $0.38\%$), so a mean difference of $0.05\%$ averages many ties ($31$--$46\%$ of the instance pairs tie even after averaging five seeds). The \emph{published} families B, C, E and F of the benchmark reviewed by \citet{OstojicEtAl2025} give disjoint sets: $525$ tuning instances, $2{,}550$ \emph{discovery} and $3{,}950$ \emph{held-out} test instances (families B, C, E), the $5{,}720$ B, C, E instances on which an offline selector is trained, and $780$ of family F, which no earlier analysis had used. Instances with times near $10^{18}$ do not fit the $64$-bit compiled core and were excluded ($378$ of the $7{,}403$ of the first three sets and $280$ of the $6{,}000$ training ones, all in family B); each exclusion is listed in the repository.

\paragraph{Unrelated machines.} In $R||\Cmax$ the processing time $p_{jk}$ of job $j$ depends on the machine $k$. We use four instance sources. (a) \emph{Synthetic}: the twelve classes above with $p_{jk}\sim U[1,100]$ independent, $100$ test and $30$ tuning instances each. (b) \emph{Published families}: the three families used in the exact-methods literature~\citep{AbladEtAl2021} (close to those of \citet{FanjulPeyroRuiz2010}): uncorrelated $p_{jk}\sim U\{10,100\}$; correlated jobs $p_{jk}=d_{jk}+b_j$ with $b_j\sim U\{1,100\}$, $d_{jk}\sim U\{1,20\}$; correlated machines $p_{jk}=a_k+d_{jk}$ with $a_k\sim U\{1,100\}$, $d_{jk}\sim U\{1,20\}$; with $m\in\{3,5,10,15,20\}$ and $n\in\{30,50,80,100,200\}$ (25 sizes per family; $75$ classes), $20$ tuning and $20$ test instances per class, generated by us from that description (the original files were unreachable when we asked). (c) The \emph{standard benchmark} of Fanjul-Peyro and Ruiz~\citep{FanjulPeyroRuiz2010}: $1{,}400$ instances ($7$ types $\times$ $m\in\{10,20,30,40,50\}\times n\in\{100,200,500,1000\}\times10$), with the best value found by CPLEX (2 h), Gurobi (300 s) and an exact Lagrangian method (300 s), run by third parties; the study uses the minimum of the three as the \emph{best known} solution, which is proven optimal for $73\%$ of the test instances. Two instances per cell tune, five test, three are reserved.

\paragraph{Flowshop and genetic algorithm.} The permutation flowshop host uses nine synthetic classes ($n\in\{20,50,100\}$, $m\in\{5,10,20\}$, $20$ tuning and $40$ test instances each) and the $90$ instances ta001--ta090 of \citet{Taillard1993}. The genetic algorithm is steady-state for $P||\Cmax$ ($1{,}200$ instances, $T=10{,}000$ offspring).

\subsection{Search mechanisms and learners}
\label{sec:learners}
\paragraph{Mechanisms.} The \textbf{IG} starts from LPT (or from the best of LPT, MULTIFIT and LDM) and alternates destruction (removing $d$ jobs) and greedy reconstruction (non-increasing time onto the machine with the earliest completion), with greedy acceptance $\Cmax(x')\le\Cmax(x)$ and no separate perturbation phase. The \textbf{ILS} starts from LPT and, at each iteration, a selector chooses one of four neighbourhood operators (inter-machine swap, reinsertion, intra-machine swap, block move) under best improvement; after $50$ iterations without improvement it restarts from the best solution. The \emph{actions} of the IG are the destruction strategy (random, longest jobs, machine-concentrated; $3$ actions), the destruction size ($d/n\in\{0.05,0.10,0.20,0.30,0.40\}$; $5$ actions), their $15$ combinations, these $15$ with or without a descent after reconstruction ($30$ actions); for the standard benchmark the size grid is scaled to $n$ ($d/n\in\{0.01,0.02,0.05,0.10,0.20\}$) because $d=50$ jobs is already too destructive for $n=1000$. The \emph{flowshop} host uses the IG of Ruiz and St\"utzle (NEH start, insertion local search, constant-temperature acceptance; the actions are the destruction size, alone or crossed with three temperature factors); the \emph{genetic algorithm} has population $20$, binary tournament, and the child replaces the worst individual (an action is a crossover crossed with a mutation and one of two intensities).

\paragraph{State, reward, selectors.} At each iteration a selector $\pi$ chooses an action from the state of the current solution (four features: gap to the bound, load imbalance, fraction of the budget used, stagnation, discretised into $81$ states), the IG destroys, reconstructs and optionally descends accordingly, $\pi$ receives the reward, the improvement of the makespan truncated at zero, and the candidate is accepted if it does not worsen the makespan; the selector is the only component that changes between arms.

\paragraph{Learners.} Tabular Q-Learning over the $81$ states ($\alpha=0.1$, $\gamma=0.9$, $\varepsilon$-greedy from $1.0$ to $0.05$ by $0.995$ per step)~\cite{WatkinsDayan1992}; UCB1~\cite{AuerEtal2002}; a policy-gradient learner (softmax preferences per state, learning rate chosen on tuning instances)~\cite{Williams1992,SuttonBarto2018}; six methods without state, all tuned on the tuning instances: probability matching, adaptive pursuit, an ALNS-style roulette, Thompson sampling with discounting~\cite{Russo2018}, DMAB~\cite{DaCostaEtal2008} and an online racing learner; LinUCB on a continuous context~\cite{LiEtAl2010}; a DQN with two layers of $64$ units, replay and a target network~\cite{MnihEtAl2015}; an optimistic Q-Learning (optimistic initial table, constant $\varepsilon=0.05$, scale-invariant reward); and two re-implemented published designs: an episodic Q-Learning with a binary state in the style of \citet{KarimiMamaghanEtal2023} and a sliding-window bandit with rank-based credit in the style of FRRMAB~\cite{LiFialhoEtal2014}; the original hyperparameters were not available, so we tuned them on the tuning instances on a small grid. \emph{Random choice} among the same actions is the control. The \emph{static competitor} is the best fixed action per class, chosen on the tuning instances.

\paragraph{Action order and ties.} The tabular Q-Learning starts from a table of zeros, its reward is non-negative, and its greedy step takes the \emph{first} maximum. In a state that has not been visited all values are equal, so after the exploration rate has decayed the learner repeats the action with the smallest index; with $81$ states, up to $30$ actions and a thousand steps most state--action pairs are never visited. Our action lists are ordered by destruction size, so the first action is the smallest and cheapest one. We therefore report the standard implementation (\textsc{ql}) \emph{and} an order-invariant version (\textsc{ql-perm}) that draws a random permutation of the actions in each run, which is equivalent to breaking ties at random; the other learners (optimistic Q-Learning, UCB1, the bandits) do not depend on the order in the same way. Section~\ref{sec:order} shows how much of the standard learner's behaviour this explains.

\subsection{The budget-fair attribution protocol}
\label{sec:protocol}
\paragraph{Terms.} A \emph{class} is an $(n,m)$ pair, a published family by size tercile or, in family F, a subset with $m$ and $n$; a \emph{cell} is a class in one action space (five or fifteen actions) and is the unit of the headroom analysis; the \emph{validated action} is the best fixed action per class, chosen on separate instances; the \emph{headroom} of a cell is the gain over random choice of the best fixed action per instance. ``Level'' means parity or no significant difference; \emph{registered} results were specified before the battery was run, \emph{post hoc} ones were not. Throughout, a \emph{gain} is $-\bar d$ in percentage points of the makespan (written ``points'' or $\%$), positive when the first-named arm is better.

\paragraph{(1) Random control at a matched budget.} The effect of the policy is the difference between the learner and uniform random choice over the same actions, iterations and instances.

\paragraph{(2) Equal CPU time, calibrated sequentially; decision and cost effects.} For each pair, the CPU time ratio $K=\mathrm{CPU}_1/\mathrm{CPU}_2$ of $1000$ iterations of the first arm to $1000$ of the second is measured on separate instances, in sequence (not always on a quiet machine), and the second arm receives $\mathrm{round}(1000K)$ iterations; since $K$ is a property of the implementation, every arm must be implemented with comparable efficiency, and $K$ is reported. Let $\bar d_{\mathrm{iter}}$ and $\bar d_{\mathrm{CPU}}$ be the mean relative differences of the learner against random choice at equal iterations and at equal CPU time. A learner's gain at equal CPU time is the sum of a \emph{decision effect}, its gain at equal iterations, and a \emph{cost effect}, the extra iterations the random control loses by matching the CPU time of a learner that chooses cheaper actions:
\begin{equation}
G_{\mathrm{CPU}}=G_{\mathrm{dec}}+G_{\mathrm{cost}},\qquad G_{\mathrm{dec}}=-\bar d_{\mathrm{iter}},\qquad G_{\mathrm{cost}}=\bar d_{\mathrm{iter}}-\bar d_{\mathrm{CPU}}.
\label{eq:decomp}
\end{equation}
The identity is an accounting one, but it separates two claims that are usually mixed: that the learner \emph{decides} better than random, and that it \emph{runs} faster by choosing cheap actions. A learner with $G_{\mathrm{dec}}\le0$ and $G_{\mathrm{cost}}>0$ wins at equal CPU time without having decided anything better.

\paragraph{(3) A static competitor validated offline.} On separate instances, the best fixed action per class and, when actions are complementary, a uniform mixture of the best few; this is what a practitioner would use. When it is compared with a learner at equal CPU time, it receives its own cost ratio.

\paragraph{(4) Cross-fitted oracle headroom.} The gain over random choice of the best fixed action per instance estimates what a selector that picks one action per instance could gain; the action is chosen on two seeds and evaluated on two fresh ones, because the best of one seed per action carries run-to-run noise (Section~\ref{sec:estimator}).

\paragraph{(5) Power and positive controls.} The minimum detectable gain is estimated analytically and by planting gains in the paired differences, and the learners and the pipeline are checked on a planted reward bonus and on a static policy chosen with class information.

\subsection{The cross-fitted headroom estimator}
\label{sec:estimator}
Let $x$ be an instance, $a\in\{1,\dots,A\}$ an action and $Y_a(x)$ the relative makespan of one run, with $\mathbb E\,Y_a(x)=\mu_a(x)$ and independent run-to-run noise of variance $\sigma^2$. Let $\mu_{\mathrm{rand}}(x)$ be the expected makespan of the random control. The \emph{headroom} of $x$ is $H(x)=\mu_{\mathrm{rand}}(x)-\min_a\mu_a(x)$ and the headroom of a cell is its mean over instances. The \emph{one-seed oracle} takes $\hat H_1=\mu_{\mathrm{rand}}-\min_aY_a$ from one run per action; the \emph{cross-fitted oracle} chooses $\hat a=\arg\min_a\bar Y_a^{(1)}$ on $k_1$ runs and evaluates it on $k_2$ independent ones, $\hat H_{\mathrm{cf}}=\mu_{\mathrm{rand}}-\bar Y_{\hat a}^{(2)}$ (we use $k_1=k_2=2$).

\begin{proposition}
\label{prop:bounds}
For every $x$, $\mathbb E[\hat H_{\mathrm{cf}}(x)]\le H(x)\le\mathbb E[\hat H_1(x)]$. If all $\mu_a(x)$ are equal, then $\mathbb E[\hat H_{\mathrm{cf}}(x)]=H(x)=0$ and $\mathbb E[\hat H_1(x)]=\sigma\,\mathbb E[\max_{a}Z_a]\approx\sigma\sqrt{2\ln A}$ for independent standard normal $Z_a$ (under normal noise).
\end{proposition}
\begin{proof}[Sketch]
The expectation of a minimum does not exceed the minimum of the expectations, so $\mathbb E[\min_aY_a]\le\min_a\mu_a$ and $\mathbb E[\hat H_1]\ge H$. Conditional on $\hat a$, which depends only on the first batch, $\mathbb E[\bar Y_{\hat a}^{(2)}\mid\hat a]=\mu_{\hat a}\ge\min_a\mu_a$, so $\mathbb E[\hat H_{\mathrm{cf}}]\le H$; the gap $\mathbb E[\mu_{\hat a}-\min_a\mu_a]$ is the expected regret of the selection, which falls as $k_1$ grows. With equal means every $\mu_{\hat a}=\mu_a$, so the regret is zero. The full argument and a Monte Carlo check are in Appendix~\ref{app:estimator}.
\end{proof}
Hence the pair $(\hat H_{\mathrm{cf}},\hat H_1)$ brackets the true headroom: the cross-fitted value is a lower bound and the one-seed value an upper bound, and the one-seed value is large even when there is nothing to gain. This is why every headroom statement in this paper uses the cross-fitted value, and why the one-seed value, which earlier versions of our analysis used, produced a spurious negative correlation with the learner's gain on the flowshop (Section~\ref{sec:rq3}).

\subsection{Statistics and registration}
\label{sec:stats}
Per instance (the mean of the seeds), $d=100(A-B)/B$ in the makespan is tested in pairs (Wilcoxon, two-sided, ties dropped), with Holm correction over each family of comparisons, a bootstrap $90\%$ interval, and an equivalence margin~\cite{Lakens2017} of $\pm0.05\%$, about an eighth of one makespan unit on the median synthetic instance; a Holm-significant difference whose interval lies inside the margin is reported as \emph{parity}, and one whose interval crosses the margin as a \emph{small win}. The headline pairs are repeated with Pratt's treatment of ties and a paired permutation test, and the registered tests are pooled under a global Benjamini--Yekutieli control (Section~\ref{sec:discussion}). Hypotheses, decision rules and probabilistic predictions were written, together with md5 hashes of the code and of the frozen inputs, in a dated research diary before each battery was run~\cite{NosekEtal2018}; each analysis script prints, for each prediction, whether it held. The diary is a private, version-controlled record; its pre-registration blocks have been compiled for deposit in a public registry.

\subsection{Implementation and reproducibility}
\label{sec:impl}
The local search was first written with Python loops, then vectorised (same moves, tie-breaking and results: $120$ of $120$ runs identical), then compiled with the IG, the destruction strategies and the agents (numba), exact on the deterministic parts and statistically equivalent on the random draws (mean differences inside $\pm0.05\%$ in five checks); it also counts evaluations. The other hosts have their own compiled cores with the same agent code. Unit tests ($1{,}549$ at the time of writing) cover the cores, the estimators, the equivalence between the compiled and the reference implementations, the seeds and the frozen choices. A set of audit scripts recomputes every number quoted in the text from the result files and fails if one does not match. All runs used one laptop (Intel Core i5-12450H, 8 cores, 16 GB, Windows 11), where $K$ was measured.
